% Options for packages loaded elsewhere
\PassOptionsToPackage{unicode}{hyperref}
\PassOptionsToPackage{hyphens}{url}
\PassOptionsToPackage{space}{xeCJK}
\documentclass[
  chinese,
  11pt,
]{ctexart}
\usepackage{xcolor}
\usepackage[margin=2.2cm]{geometry}
\usepackage{amsmath,amssymb}
\setcounter{secnumdepth}{5}
\usepackage{iftex}
\ifPDFTeX
  \usepackage[T1]{fontenc}
  \usepackage[utf8]{inputenc}
  \usepackage{textcomp} % provide euro and other symbols
\else % if luatex or xetex
  \usepackage{unicode-math} % this also loads fontspec
  \defaultfontfeatures{Scale=MatchLowercase}
  \defaultfontfeatures[\rmfamily]{Ligatures=TeX,Scale=1}
\fi
\usepackage{lmodern}
\ifPDFTeX\else
  % xetex/luatex font selection
  \ifXeTeX
    \usepackage{xeCJK}
    \setCJKmainfont[]{Noto Serif SC}
    \setCJKsansfont[]{Noto Sans SC}
    \setCJKmonofont[]{DejaVu Sans Mono}
  \fi
  \ifLuaTeX
    \usepackage[]{luatexja-fontspec}
    \setmainjfont[]{Noto Serif SC}
  \fi
\fi
% Use upquote if available, for straight quotes in verbatim environments
\IfFileExists{upquote.sty}{\usepackage{upquote}}{}
\IfFileExists{microtype.sty}{% use microtype if available
  \usepackage[]{microtype}
  \UseMicrotypeSet[protrusion]{basicmath} % disable protrusion for tt fonts
}{}
\makeatletter
\@ifundefined{KOMAClassName}{% if non-KOMA class
  \IfFileExists{parskip.sty}{%
    \usepackage{parskip}
  }{% else
    \setlength{\parindent}{0pt}
    \setlength{\parskip}{6pt plus 2pt minus 1pt}}
}{% if KOMA class
  \KOMAoptions{parskip=half}}
\makeatother
\usepackage{longtable,booktabs,array}
\newcounter{none} % for unnumbered tables
\usepackage{calc} % for calculating minipage widths
% Correct order of tables after \paragraph or \subparagraph
\usepackage{etoolbox}
\makeatletter
\patchcmd\longtable{\par}{\if@noskipsec\mbox{}\fi\par}{}{}
\makeatother
% Allow footnotes in longtable head/foot
\IfFileExists{footnotehyper.sty}{\usepackage{footnotehyper}}{\usepackage{footnote}}
\makesavenoteenv{longtable}
\ifLuaTeX
\usepackage[bidi=basic,shorthands=off,provide=*]{babel}
\else
\usepackage[bidi=default,shorthands=off,provide=*]{babel}
\fi
\ifLuaTeX
  \usepackage{selnolig} % disable illegal ligatures
\fi
\setlength{\emergencystretch}{3em} % prevent overfull lines
\providecommand{\tightlist}{%
  \setlength{\itemsep}{0pt}\setlength{\parskip}{0pt}}
\usepackage{bookmark}
\IfFileExists{xurl.sty}{\usepackage{xurl}}{} % add URL line breaks if available
\urlstyle{same}
\hypersetup{
  pdftitle={课题09-可验证环境},
  pdflang={zh-CN},
  hidelinks,
  pdfcreator={LaTeX via pandoc}}

\title{课题09-可验证环境}
\author{}
\date{}

\begin{document}
\maketitle

{
\setcounter{tocdepth}{2}
\tableofcontents
}
\section{课题09 开题简报 · 可验证环境 + harness}\label{ux8bfeux989809-ux5f00ux9898ux7b80ux62a5-ux53efux9a8cux8bc1ux73afux5883-harness}

\begin{quote}
模块：\textbf{M7/M8} ｜ 周次：\textbf{W13（12-07→12-13）} ｜ 算力：\textbf{T2 Kaggle 2×T4，计划 6 GPU·小时} 唐杰作业对应：\textbf{⑤ 搭可验证环境 + harness} ｜ 判分来源：\textbf{自建}（判定逻辑单测 + rollout 成功率） 手写：\textbf{R9（判定逻辑 / verifier）}
\end{quote}

\subsection{一、研究背景}\label{ux4e00ux7814ux7a76ux80ccux666f}

2026 年的关键词是''\textbf{可验证奖励}''：与其训练一个会拍马屁的奖励模型，不如\textbf{造一个能判定对错的环境}。 CS336 没有这一环，Berkeley Agentic AI MOOC 专门有一讲叫 ``Post-Training Verifiable Agents''， \texttt{PrimeIntellect-ai/verifiers} 把它抽象成库。唐杰作业⑤要求''搭可验证环境 + harness''正是这个方向。

\textbf{``可验证''的三要素}：①任务是可判定的（不是主观好坏）②环境能给出稳定反馈 ③harness 能重复运行并复现。

\subsection{二、研究问题}\label{ux4e8cux7814ux7a76ux95eeux9898}

\begin{quote}
\textbf{我能为一个小模型设计一个''可自动判定''的任务环境吗？判定逻辑的假阳性/假阴性有多高？模型在这个环境里的成功率是多少？}
\end{quote}

\subsection{三、攻关线索}\label{ux4e09ux653bux5173ux7ebfux7d22}

\begin{enumerate}
\def\labelenumi{\arabic{enumi}.}
\tightlist
\item
  \textbf{选题}：从三个里选一个（按成本从低到高）

  \begin{itemize}
  \tightlist
  \item
    ①\textbf{数学题解答}（用数值比对 + 容错；MATH-500 风格）
  \item
    ②\textbf{函数补全}（隐藏单元测试；本质是迷你 HumanEval）
  \item
    ③\textbf{格式约束生成}（如严格 JSON / 表格填充，规则判定）
  \end{itemize}
\item
  \textbf{判定逻辑的三大坑}：等价但不相同（\texttt{1/2} vs \texttt{0.5}）、解析失败当错、\textbf{边界用例漏判}
\item
  \textbf{假阳性检测}：人为构造''错误但能骗过判定器''的答案 → 测出\textbf{假阳性率}（这是质量的核心指标）
\item
  \textbf{harness 设计}：任务装载 → 模型生成 → 判定 → 记录轨迹 → 复现（固定种子）
\item
  \textbf{可重复性}：同一批任务跑两次，成功率波动多大？（\textbf{误差棒思维}，M7 讲过）
\item
  \textbf{对照}：可验证奖励 vs 规则弱判定（如只看格式），用同一模型比成功率差异
\end{enumerate}

\subsection{四、里程碑}\label{ux56dbux91ccux7a0bux7891}

\begin{itemize}
\tightlist
\item[$\square$]
  M1 选定任务类型，写 30--100 道题的\textbf{题集}（含标准答案与判定规则）
\item[$\square$]
  M2 \textbf{判定逻辑单测}：等价形式、边界、异常输入（≥10 个用例）
\item[$\square$]
  M3 \textbf{假阳性实验}：构造 10 个''错误但可能蒙混''的答案，测出漏判率
\item[$\square$]
  M4 harness 跑通：给基座模型跑一遍，得成功率 + 失败样本
\item[$\square$]
  M5 失败分类：把失败归成 3--5 类（格式错 / 计算错 / 中途放弃 / 判定器误杀）
\item[$\square$]
  M6 一页报告：\textbf{判定器质量 + 模型表现 两个都要报}
\end{itemize}

\subsection{五、交付物}\label{ux4e94ux4ea4ux4ed8ux7269}

{\def\LTcaptype{none} % do not increment counter
\begin{longtable}[]{@{}lll@{}}
\toprule\noalign{}
交付物 & 位置 & 验收 \\
\midrule\noalign{}
\endhead
\bottomrule\noalign{}
\endlastfoot
题集与判定器 & \texttt{手写/}（判定逻辑 R9） & M2 单测通过 \\
harness 脚本 & \texttt{脚手架/}（B3） & 固定种子可复跑 \\
失败样本 & \texttt{证据/} & ≥20 条，分类标注 \\
假阳性报告 & \texttt{报告.md} & 含漏判率与典型案例 \\
\end{longtable}
}

\subsection{六、讲课锚点}\label{ux516dux8bb2ux8bfeux951aux70b9}

\begin{itemize}
\tightlist
\item
  \textbf{讲义}：\texttt{M8-智能体.md}（W13 展开）
\item
  \textbf{对标}：Berkeley Agentic AI MOOC 的 ``Post-Training Verifiable Agents'' 与 ``Agent Evaluation'' 讲 · \texttt{PrimeIntellect-ai/verifiers}（\textbf{只读概念，不引入代码}）· \texttt{harbor-framework/terminal-bench} 的任务设计思路 · \textbf{rasbt ch3}（评测器设计：解析 + verifier + 基准式检查）
\end{itemize}

\subsection{七、代码级提问示例}\label{ux4e03ux4ee3ux7801ux7ea7ux63d0ux95eeux793aux4f8b}

\begin{enumerate}
\def\labelenumi{\arabic{enumi}.}
\tightlist
\item
  你的判定器把 \texttt{1/2} 判错、把 \texttt{0.50} 判对------这属于哪一类错误？如何系统性避免？
\item
  假阳性率高会导致 RL 训练出现什么后果？（联系课题08 的 reward hacking）
\item
  同一批任务两次运行成功率差 8\%。你会先怀疑模型还是 harness？怎么验证？
\item
  为什么''用 LLM 当裁判''在可验证任务上通常是\textbf{退步}？
\end{enumerate}

\subsection{八、风险与提醒}\label{ux516bux98ceux9669ux4e0eux63d0ux9192}

\begin{itemize}
\tightlist
\item
  ⚠️ \textbf{判定器本身也要被评测}：只报模型成功率不报判定器质量，是不完整的实验。
\item
  ⚠️ 题集不要抄公开 benchmark（污染问题），\textbf{自己出题 + 自己判定}才是本课题的价值。
\item
  ⚠️ 配额有限：优先''题集质量高、任务简单''，不要一上来做代码执行沙箱（那是课题10 的事）。
\end{itemize}

\end{document}
